\subsection{Excitation and readout, case by case}
\label{si:whatsees}

Every channel in the atlas is specified by two independent choices: what
the pulse \emph{drives}, fixed by the field orientation, and what the
analyser \emph{reads}, fixed by which Cartesian component of the total
magnetic moment radiates into the accepted polarisation.  For
propagation along $b$ the two transverse polarisation states are
($H\!\parallel\!a$, $E\!\parallel\!c$) and ($H\!\parallel\!c$,
$E\!\parallel\!a$), so ``same'' and ``cross'' detection select $m_x$ and
$m_z$ respectively in geometry A, and the reverse in geometry B.

\begin{table}[h]
\centering\small
\caption{Excitation and readout for the four measured configurations.
$d_z^{\rm eff}$ is the exchange-induced $c$-axis dipole of the $E_{12}$
transition; $\mu^z_{23}=0.913$ and $\mu^z_{13}\simeq0$ (the selection
rule of \cref{si:peaksAsame}); $\mu^z_{12}=5.264$ is the Tm $c$-axis
magnetic moment; $\beta\simeq7\mu$ is the hyperpolarisability allowed by
the absence of inversion at the Tm $4c$ site.  Fe$^{3+}$ at $4b$ sits on
an inversion centre and therefore contributes no quadratic term.}
\label{tab:si-channels}
\renewcommand{\arraystretch}{1.35}
\begin{tabular}{p{2.1cm}p{5.2cm}p{7.4cm}}
\toprule
configuration & excitation (what the pulse drives) & readout (what the
analyser detects)\\
\midrule
\textbf{A, cross}\newline $H\!\parallel\!a$, $E\!\parallel\!c$ in;
$H\!\parallel\!c$ out
& Fe: Zeeman $H_x F_x$ $\Rightarrow$ $\qafm$ (which carries
  $\delta F_x$).\newline
  Tm: $E\!\parallel\!c$ on the $c$-axis dipoles, drive vector
  $(\lambda^2,\lambda^5,\lambda^7)\!\propto\!(d_z^{\rm eff},0,\mu^z_{23})$
  $\Rightarrow$ $E_{12}$ and $E_{23}$; $E_{13}$ \emph{dark}.
& $m_z=F_z+\mu^z_{12}\lambda^2+\beta\lambda^1\lambda^2$.\newline
  $F_z\equiv0$ (machine zero: $\qafm$ has no $c$-axis moment), so the
  channel is \emph{purely} Tm: the $E_{12}$ dipole, its $W^{xz}_1$-forced
  $\qafm$ sideband, and the quadratic term radiating at $2E_{12}$.\\
\textbf{A, same}\newline $H\!\parallel\!a$, $E\!\parallel\!c$ in;
$H\!\parallel\!a$ out
& identical to A-cross --- it is the same experiment, differing only in
  the analyser.
& $m_x=F_x+2.39\lambda^5+0.91\lambda^7$ (quadratures
  $\lambda^{4},\lambda^{6}$).\newline
  Fe $F_x$ carries $\qafm$; the Tm $a$-axis moments carry $E_{13}$ and
  $E_{23}$.  This is where the CEF--CEF transfers appear.\\
\textbf{B, cross}\newline $H\!\parallel\!c$, $E\!\parallel\!a$ in;
$H\!\parallel\!a$ out
& Fe: Zeeman $H_z F_z$ $\Rightarrow$ $\qfm$ (which carries
  $\delta F_z$, $\delta F_y$).\newline
  Tm: the $c$-directed field projects onto $\mu^z_{12}$ only, driving
  $\lambda^2$ alone $\Rightarrow$ the $\{\lambda^1,\lambda^2,\lambda^3\}$
  subalgebra closes.
& $m_x=F_x+2.39\lambda^5+0.91\lambda^7$.\newline
  The Tm part is \emph{identically zero} ($\lambda^{4\text{--}7}\equiv0$
  by the closure above), so this channel is pure Fe: a magnon map read
  out on the $\qafm$ line.\\
\textbf{B, same}\newline $H\!\parallel\!c$, $E\!\parallel\!a$ in;
$H\!\parallel\!c$ out
& identical to B-cross.
& $m_z=F_z+\mu^z_{12}\lambda^2+\beta\lambda^1\lambda^2$.\newline
  $F_z$ carries $\qfm$, but the Tm $\lambda^2$ is driven directly here
  and is $50\times$ larger, so the channel is Tm-dominated --- see the
  prediction below.\\
\bottomrule
\end{tabular}
\end{table}

Two entries are exact statements rather than approximations, and both do
real explanatory work.  First, $F_z\equiv0$ in geometry A: the $\qafm$
magnon lives in $\delta F_x$ and has no $c$-axis moment, so the Fe
sublattice is radiatively silent in the A-cross channel and anything
seen there at a magnon frequency must come from the Tm ion.  Second, in
geometry B the $c$-directed drive couples only to $\mu^z_{12}$, closing
the $\{\lambda^1,\lambda^2,\lambda^3\}$ subalgebra, so
$\lambda^{4\text{--}7}$ vanish identically and no $E_{13}$ or $E_{23}$
emission is possible in either geometry-B channel.
